\documentclass{article}
%%%%% PROGRAMMING PACKAGES
\usepackage{expl3,xparse}     %LaTeX3
\usepackage{mfirstuc}    %first letter uppercase with \makefirstuc (used in todo notes)
\usepackage{settobox}     %get box dimensions with \settoboxwidth (used in fit environment) => pretty stupid package.. could be avoided

%%%%% FIGURES, BOXES, COLORS
\usepackage{graphics,graphicx,trimclip,adjustbox}     %figures and box management
\usepackage{color,xcolor}     %to use colors (used in the todonotes)
    
%%%%% ENUMERATE
\usepackage{enumitem}   %to use "enumerate"
    \setlist[enumerate,1]{label = $\bullet$}   %default label at bullet points
    
%%%%% TODO NOTES
\usepackage[textwidth=16.3mm,textsize=scriptsize]{todonotes}   %to put notes everywhere
\newcommand*{\newcommentator}[2][red!20]{%
    \expandafter\newcommand\csname #2\endcsname[1]{\todo[color=#1]{{\bf\makefirstuc{#2}:} ##1}}%
    \expandafter\newcommand\csname big#2\endcsname[1]{\todo[inline,color=#1]{{\bf\makefirstuc{#2}:} ##1}}}
\newcommentator{pablo}

%%%%% MY HYPERLINK AND REF/LABEL MACROS
\usepackage{mypreamble/linkref}

%%%%% MATH PACKAGES
\usepackage{amssymb,amsmath}     %math original package
    \numberwithin{equation}{section}
\usepackage{mathtools}     %improves AMSmath
\usepackage{stackrel}     %better stacking of math symbols
\usepackage{mleftright}     %defines \mleft and \mright
    \mleftright     %redefines \left as \mleft and \right as \mright
\usepackage{pgfcore}
\usepackage{tikz-cd}     %advanced and rigorous diagrams
    
%%%%% MY SPECIAL PACKAGES
\usepackage{mypreamble/mathenv}     %improves on AMSmath's equation
\usepackage{mypreamble/thmenv}     %replaces the AMSthm package

%%%%% MATH MACROS AND FONT PACKAGES
%\usepackage{fontspec}
\usepackage{mypreamble/mathcro}

%%%%% HYPERREF PACKAGE
\makeatletter
\pretocmd\refstepcounter{%
 \zref@setcurrent{counter}{#1}}{}{}
\makeatother
\usepackage[pdfencoding=auto]{hyperref}   %hyperlinks parts of the pdf
    \hypersetup{
        bookmarksopen=true,   %expends tree of bookmarks in Acrobat pdf reader
        bookmarksnumbered=true,  %show numbers of sections in bookmarks
        linktoc=all,   %both section and page numbers are links
        hidelinks   %links (color+border) are hidden
    }
\providecommand*{\definitionautorefname}{Definition}
\providecommand*{\notationautorefname}{Notation}
\providecommand*{\conventionautorefname}{Convention}
\providecommand*{\conjectureautorefname}{Conjecture}
\providecommand*{\assumptionautorefname}{Assumption}
\providecommand*{\remarkautorefname}{Remark}
\providecommand*{\observationautorefname}{Observation}
\providecommand*{\lemmaautorefname}{Lemma}
\providecommand*{\propositionautorefname}{Proposition}
\providecommand*{\corollaryautorefname}{Corollary}
\providecommand*{\exampleautorefname}{Example}
\providecommand*{\exerciseautorefname}{Exercise}
\providecommand*{\openquestionautorefname}{Open Question}

%%%%% OUTPUT LAYOUT PACKAGES
\usepackage{geometry}     %to set margins
\geometry{
    a4paper,
    total={170mm,257mm},
    left=20mm,
    top=20mm,
    marginparwidth=20mm,
}
%% Clearpage before sections but not before TOC
\let\oldsection\section
\renewcommand\section{\clearpage\oldsection}
\makeatletter
\renewcommand\tableofcontents{%
    \oldsection*{\contentsname\@mkboth{\MakeUppercase\contentsname}{\MakeUppercase\contentsname}}\@starttoc{toc}%
}
\makeatother    %ensures the TOC stays in front page

%%%%% TYPESET DEBUGGING
%\usepackage{showframe}     %to make margins appear
%\usepackage{lua-visual-debug}     %to make all boxes appear
%\overfullrule=1pt   %creates 1pt wide black bars next to any overfull hbox

%%%%% QUIVER SETUP
\usetikzlibrary{calc,decorations.pathmorphing}
% A TikZ style for curved arrows of a fixed height, due to AndréC.
\tikzset{curve/.style={settings={#1},to path={(\tikztostart)
    .. controls ($(\tikztostart)!\pv{pos}!(\tikztotarget)!\pv{height}!270:(\tikztotarget)$)
    and ($(\tikztostart)!1-\pv{pos}!(\tikztotarget)!\pv{height}!270:(\tikztotarget)$)
    .. (\tikztotarget)\tikztonodes}},
    settings/.code={\tikzset{quiver/.cd,#1}
        \def\pv##1{\pgfkeysvalueof{/tikz/quiver/##1}}},
    quiver/.cd,pos/.initial=0.35,height/.initial=0}

% TikZ arrowhead/tail styles.
\tikzset{tail reversed/.code={\pgfsetarrowsstart{tikzcd to}}}
\tikzset{2tail/.code={\pgfsetarrowsstart{Implies[reversed]}}}
\tikzset{2tail reversed/.code={\pgfsetarrowsstart{Implies}}}
% TikZ arrow styles.
\tikzset{no body/.style={/tikz/dash pattern=on 0 off 1mm}}
\usepackage{mathdots}
\usepackage{bookmark}

%%%%% TITLE PAGE
\title{Solutions 321-1}
\author{Pablo Bustillo Vazquez}

\begin{document}

\maketitle
\tableofcontents
\newpage

\section{Week 1}

\begin{example}[Exercice 2] We have to prove that $\alpha = \sup(E)$ if and only if:
\[\begin{cases}[l]
    \forall x \inset E, \quad x \leq \alpha \\
    \forall \epsilon > 0, \quad \exists x \inset E, \quad \alpha - \epsilon \leq x
\end{cases}\]
The first condition is precisely the definition of upper bound. The second condition can be rewritten as
\begin{equation*}[c+l+l]
    & \forall \beta, \quad \beta < \alpha \quad \implies \quad \exists x \inset E, \quad \beta \leq x &\\
    \iff    & \forall \beta, \quad \beta < \alpha \quad \implies \quad \beta \textnormal{ is not an upper bound} & \textnormal{(by defintion)}\\
    \iff    & \forall \beta, \quad \beta \textnormal{ is an upper bound} \quad \implies \quad \alpha \leq \beta & \textnormal{(by contraposition)}\\
\end{equation*}
which adds that $\alpha$ is the smallest upper bound.
\end{example}

\begin{example}[Exercice 3]\leavevmode
\begin{enumerate}[label = (\alph*)]
    \item The set $\mathbb{Z}_{\leq x} = \{n \inset \mathbb{Z} \mid n \leq x\}$ is bounded above by $x$ and contains $0$, hence has a supremum that we will denote $\lfloor x \rfloor$.
    Because $x$ is an upper bound, we clearly have $\lfloor x \rfloor \leq x$. By exerice 2, with $\epsilon = 1$, there exists $n \leq x$ such that
    \begin{equation}[c+l]
        \nonumber & \lfloor x \rfloor - 1 < n \leq \lfloor x \rfloor \\
       \label{first_ineq}\implies & \lfloor x \rfloor  < n + 1
    \end{equation}
    Assume for the sake of contradiction that $\epsilon = \lfloor x \rfloor - n > 0$. By exercice 2 again, we have $m \leq x$ such that
    \begin{equation*}[c+l+l]
        & n = \lfloor x \rfloor - \epsilon < m \leq \lfloor x \rfloor &\\
       \implies & \lfloor x \rfloor < n + 1 \leq m \leq \lfloor x \rfloor & \textnormal{(by \autoref{first_ineq} and by successor property)}
    \end{equation*} 
    which is a contradiction. Hence $\lfloor x \rfloor = n$ is an integer and hence a maximum of $\mathbb{Z}_{\leq x}$.

    \item We have 
    \begin{equation*}[c+l]
        & \sum_{j=0}^{k-1} \frac{n_j}{10^j} \leq x \\
       \iff &  n_k \leq x 10^k - \sum_{j=0}^{k-1} n_j 10^{k-j}
    \end{equation*}
    Hence the sets of all $n_k \inset \mathbb{Z}$ satisfying either property are identical. By part $(a)$,
    \[
        n_k \coloneqq \left\lfloor x 10^k - \sum_{j=0}^{k-1} n_j 10^{k-j} \right\rfloor
    \]
    satisfies the desired property.

    \item Recall that in part $(a)$, we established that: for all $x \inset \mathbb{R}_{\geq 0}$, we have $\lfloor x \rfloor \inset \mathbb{Z}_{\geq 0}$ such that $\lfloor x \rfloor \leq x < \lfloor x \rfloor + 1$.
    In particular, using the construction of part $(b)$ inductively, we automatically have
    \begin{equation}[c+l]
        \nonumber & n_k \leq x 10^k - \sum_{j=0}^{k-1} n_j 10^{k-j} < n_k + 1 \\
        \label{decimal} \iff & 0 \leq x - \sum_{j=0}^{k} \frac{n_j}{10^j} < \frac{1}{10^k}
    \end{equation}
    for all $k \geq 0$. Now for the second part, take $k \geq 1$, using \autoref{decimal} with $k-1$ we get
    \begin{equation*}[c+l]
        & 0 \leq x 10^k - \sum_{j=0}^{k-1} n_j 10^{k-j} < 10 \\
        \implies & 0 \leq n_k \leq x 10^k - \sum_{j=0}^{k-1} n_j 10^{k-j} < 10
    \end{equation*}
    as required, where the first inequality comes from $0$ being in the defining set whose supremum is $n_k$.

    \item For all $x > 0$, 
    \[
        (1 + x)^k = 1 + kx + \sum_{i=2}^{n} \binom{k}{i} x^i \geq 1 + kx
    \]
    Taking $x  = c-1$ in the above, if $\{c^k \mid k \inset \mathbb{N}\}$ is bounded above by $M$, then $\mathbb{N}$ is bounded above by $\frac{M-1}{c - 1}$.
    Using part $(a)$, this makes it impossible for the successor integer $\left \lfloor \frac{M-1}{c - 1} \right \rfloor + 1$ to exist.
    By contraposition, $\{c^k \mid k \inset \mathbb{N}\}$ must be unbounded.

    \item It's clear that this set is nonempty and bounded bellow by $0$.
    Using exercice $2$, we show that $0$ is the infimum.
    Let $\epsilon > 0$.
    By part $(d)$, we know that $1/\epsilon$ is not an upper bound of $\{c^k \mid k \inset \mathbb{N}\}$.
    Hence there exists $k \inset \mathbb{N}$ such that 
    \begin{equation*}[c+l]
        & 1/\epsilon < c^k \\
        \implies & 0 \leq c^{-k} < \epsilon
    \end{equation*}
    proving that
    \[
        0 = \inf \{c^{-k} \mid k \inset \mathbb{N}\}
    \]

    \item This set is non-empty by construction.
    By definition of the $n_k$'s, we clearly have that $x$ is an upper bound for this set.
    Using exercice $2$, we show that $x$ is the supremum.
    Let $\epsilon > 0$.
    Applying part $(e)$ to $c = 10$, there exists $k \geq 0$ such that $10^{-k} < \epsilon$.
    Hence by \autoref{decimal}
    \[
        x - \epsilon < x - \frac{1}{10^k} < \sum_{j=0}^{k} \frac{n_j}{10^j} \leq x
    \]
    showing that $x$ is the supremum.

    \item This set is non-empty by construction. For all $k \geq 0$,
    \[
        \sum_{j=0}^{k} \frac{m_j}{10^j} \leq m_0 + \sum_{j=1}^{k} \frac{9}{10^j} = m_0 + 1 - \frac{1}{10^k} < m_0 + 1
    \]
    Hence $m_0 + 1$ is an upper bound for this set.

    \item Let $A = \{ k \inset \mathbb{N} \mid n_k \neq m_k\}$.
    Assume for the sake of contradiction that this set is non-empty.
    By the well-ordering principle, $A$ has a minimum $k_0$.
    Without loss of generality, assume $n_{k_0} < m_{k_0}$.
    By assumption we can find $k_1 \geq k_0$ such that either $n_{k_1} \neq 9$.
    By exercice $2$ with $\epsilon = 10^{-k_1}$, we have $k_2$ and $k_3$ such that
    \[\begin{cases}
            x - \frac{1}{10^{k_1}} < \sum_{j=0}^{k_3} \frac{n_j}{10^j}\\
            x - \frac{1}{10^{k_1}} < \sum_{j=0}^{k_4} \frac{m_j}{10^j}
    \end{cases}\]
    Let $k = \max(k_1, k_2, k_3)$.
    Because the sequences considered are increasing, we have
    \begin{equation*}[l+c]
        & \begin{cases}
            x - \frac{1}{10^k} < \sum_{j=0}^{k} \frac{n_j}{10^j} < x\\
            x - \frac{1}{10^k} < \sum_{j=0}^{k} \frac{m_j}{10^j} < x
        \end{cases}\\
        \implies & \sum_{j = k_0}^{k} \frac{m_j - n_j}{10^j} =  \sum_{j = 0}^{k} \frac{m_j - n_j}{10^j} < \frac{1}{10^{k_1}}
    \end{equation*}
    However, we also have
    \begin{equation*}[rCl]
        \sum_{j = k_0}^{k} \frac{m_j - n_j}{10^j} &  \geq & \frac{m_{k_0} - n_{k_0}}{10^{k_0}} - \sum_{j = k_0 + 1}^{k_1 - 1} \frac{|m_j - n_j|}{10^j} + \frac{m_{k_1} - n_{k_1}}{10^{k_1}} - \sum_{j = k_1 + 1}^{k} \frac{|m_j - n_j|}{10^j} \\
        &\geq & \frac{1}{10^{k_0}} - \sum_{j = k_0 + 1}^{k_1 - 1} \frac{9}{10^j} - \frac{8}{10^{k_1}} - \sum_{j = k_1 + 1}^{k} \frac{9}{10^j} \\
        & = & \frac{1}{10^{k_1}} + \frac{1}{10^{k}}\\
        & > & \frac{1}{10^{k_1}}
    \end{equation*}
    which is contradiction.
    Therefore $A = \emptyset$ and we have as required: for all $k \geq 0$, $n_k = m_k$.
\end{enumerate}
    
\end{example}

\end{document}